\BourbakiExerciseGroup

\begin{enumerate}
\def\labelenumi{\arabic{enumi})}
\tightlist
\item
  \begin{enumerate}
  \def\labelenumii{\alph{enumii})}
  \tightlist
  \item
    设 \(q\) 是某个素数的幂。证明：对于每个整数 \(n > 1\) ，若 \(h_i\) （ \(1 \leqslant i \leqslant r\) ）是整数 \(n\) 的全部不同素因子，则满足 \(\zeta \in \mathbf{F}_{q^n}\) 且 \(\mathbf{F}_{q^n} = \mathbf{F}_q(\zeta)\) 的元素的个数等于
  \end{enumerate}
\end{enumerate}

\[
\nu = q ^ {n} - \sum_ {i} q ^ {n / h _ {i}} + \sum_ {i <   j} q ^ {n / h _ {i} h _ {j}} - \sum_ {i <   j <   k} q ^ {n / h _ {i} h _ {j} h _ {k}} + \dots + (- 1) ^ {r} q ^ {n / h _ {1} h _ {2} \dots h _ {r}}
\]

（注意，这样的元素由不属于任何一个域 \(F_{q}^{n/h_{i}}\) 的性质来刻画）。若 \(h_{1} \leqslant \cdots \leqslant h_{r}\) ，我们有

\[
q ^ {n} - \sum_ {i} q ^ {n / h _ {i}} \leqslant \nu \leqslant q ^ {n} - q ^ {n / h _ {1}}.
\]

考察 \(n\) 是素数之幂的情形。

\begin{enumerate}
\def\labelenumi{\alph{enumi})}
\setcounter{enumi}{1}
\item
  设 \(b_{n}(q)\) 为次数为 \(n\) 、位于 \(\mathbf{F}_q\) 上的首一不可约多项式的个数。试由 \(a\) ) 导出 \(b_{n}(q)\) 的一个表达式。特别地，若 \(q \geqslant 3\) ，证明次数为 \(n\) 、属于 \(\mathbf{F}_q[\mathbf{X}]\) 的不可约多项式（未必首一）的个数至少等于 \(q^n (q - 2) / n \geqslant q^{n + 1} / 3n\) 。
\item
  证明 \(\sum_{d\mid n}db_d(q) = q^n\) （建立恒等式 \(1 / (1 - q\mathrm{T}) = \prod_{n}(1 - \mathrm{T}^n)^{-b_n(q)}\) ）。
\end{enumerate}

* 利用默比乌斯反演公式（Lie, II, 第 176 页，附录），证明

\[
n b _ {n} (q) = \sum_ {d \mid n} \mu (n / d) q ^ {d}.
\]

将这一结果与 \(b\) 中所得的结果比较。 *

\begin{enumerate}
\def\labelenumi{\alph{enumi})}
\setcounter{enumi}{3}
\tightlist
\item
  设 \(a_{n}(q)\) 为 \(\mathbf{F}_q[\mathbf{X}]\) 中无重因式的首一多项式的个数。证明 \(a_0(q) = 1, \quad a_1(q) = q, \quad a_n(q) = q^n - q^{n-1}\) 对 \(n > 1\) 成立（建立恒等式 \(q^n = \sum_i a_{n-2i}(q) q^i\) ）。
\end{enumerate}

\begin{enumerate}
\def\labelenumi{\arabic{enumi})}
\setcounter{enumi}{1}
\item
  证明次数为 \(n\) 、属于 \(\mathbf{F}_q[\mathbf{X}]\) 且在 \(\mathbf{F}_q\) 中没有根的首一多项式的个数是 \(q^{n - q}(q - 1)^q\) ，其中 \(n \geqslant q\) （注意每个函数 \(\mathbf{F}_q \to \mathbf{F}_q\) 都是多项式函数）。
\item
  设 \(\mathbf{K}\) 是含 \(q\) 个元素的有限域；对每个素数 \(l\) ，令 \(N_{l}\) 为 \(\mathbf{K}\) 的满足下列条件的扩张的并：它们包含在 \(\Omega\) —— \(\mathbf{K}\) 的一个代数闭包 —— 之中，且次数均为 \(l\) 的幂。证明 \(\mathbf{N}_{l}\) 是 \(\mathbf{K}\) 的一个阿贝尔扩张，其伽罗瓦群同构于群 \(\mathbf{Z}_{l}\) ，即 \(l\) 进整数所成的群。由此推出 \(\Omega\) 同构于张量积 \(\otimes \mathbf{N}_{l}\) ，从而
\end{enumerate}

\(\operatorname{Gal}(\Omega / \mathbf{K})\) 同构于乘积 \(\prod_{l} \mathbf{Z}_l\) 。

* 4) 将多项式 \(\mathbf{X}^4 + 1 \in \mathbf{F}_p[\mathbf{X}]\) 分解为不可约因式。 *

\begin{enumerate}
\def\labelenumi{\arabic{enumi})}
\setcounter{enumi}{4}
\tightlist
\item
  设 \(\mathbf{K}\) 是含 \(q\) 个元素的有限域，其中 \(q\) 不是 2 的幂，并设 \(a_1, a_2, b\) 是 \(\mathbf{K}\) 的三个元素，使得 \(a_1a_2 \neq 0\) 。证明：数目 \(\nu\) （即解 \((x_1, x_2) \in \mathbf{K}^2\) 的数目，这些解满足方程 \(a_1x_1^2 + a_2x_2^2 = b\) ）由下列公式给出：
\end{enumerate}

\(1^{\circ}\) 若 \(b = 0\) 且 \(-a_{1}a_{2}\) 在 \(\mathbf{K},\nu = 1\)

\(2^{\circ}\) 若 \(b \neq 0\) 且 \(-a_{1}a_{2}\) 在 \(\mathbf{K}\) 中不是平方元，则 \(\nu = q + 1\) ；

\(3^{\circ}\) 若 \(b = 0\) 且 \(-a_{1}a_{2}\) 是 \(\mathbf{K}\) 中的平方元，则 \(\nu = 2q - 1\) ；

\(4^{\circ}\) 若 \(b \neq 0\) 且 \(-a_{1}a_{2}\) 是 K 中的平方元，则 \(\nu = q - 1\) 。

（当 \(-a_{1}a_{2}\) 不是平方元时，向 K 添加 \(\mathbf{X}^{2} + a_{1}a_{2}\) 的一个根，并利用有限域的乘法群是循环群这一事实。）

\begin{enumerate}
\def\labelenumi{\arabic{enumi})}
\setcounter{enumi}{5}
\tightlist
\item
  设 \(\mathbf{K}\) 是含 \(q\) 个元素的有限域，其中 \(q\) 不是 2 的幂。
\end{enumerate}

\begin{enumerate}
\def\labelenumi{\alph{enumi})}
\tightlist
\item
  设 \(a_{1}, a_{2}, \ldots, a_{2m}, b\) 是 K 的元素，使得 \(a_{1}a_{2} \ldots a_{2m} \neq 0\) 。证明：数目 \(\nu\) （即解 \((x_{1}, \ldots, x_{2m}) \in \mathbf{K}^{2m}\) 的数目，这些解满足方程 \(a_{1}x_{1}^{2} + \cdots + a_{2m}x_{2m}^{2} = b\) ）由下列公式给出：
\end{enumerate}

\(1^{\circ}\) 若 \(b = 0\) 且 \((-1)^{m}a_{1}a_{2}\ldots a_{2m}\) 在 K 中不是平方元，则 \(\nu = q^{2m - 1} - q^{m} + q^{m - 1}\) ；

\(2^{\circ}\) 若 \(b \neq 0\) 且 \((-1)^{m}a_{1}a_{2} \ldots a_{2m}\) 在 K 中不是平方元，则 \(\nu = q^{2m-1} + q^{m-1}\) ；

\(3^{\circ}\) 若 \(b = 0\) 且 \((-1)^{m}a_{1}a_{2}\ldots a_{2m}\) 是 K 中的平方元，则 \(\nu = q^{2m - 1} + q^{m} - q^{m - 1}\) ；

\(4^{\circ}\) 若 \(b \neq 0\) 且 \((-1)^{m}a_{1}a_{2} \ldots a_{2m}\) 是 K 中的平方元，则 \(\nu = q^{2m - 1} - q^{m - 1}\) 。

（对 \(m\) 作归纳论证，并利用习题 5。）

\begin{enumerate}
\def\labelenumi{\alph{enumi})}
\setcounter{enumi}{1}
\tightlist
\item
  设 \(a_1, a_2, \ldots, a_{2m+1}\) ，\(b\) 为 \(K\) 的元素，使得 \(a_1a_2 \ldots a_{2m+1} \neq 0\) 。证明：数目 \(\nu\) ，即下方方程的解 \((x_1, x_2, \ldots, x_{2m+1}) \in K^{2m+1}\) 的数目，
\end{enumerate}

\[
a _ {1} x _ {1} ^ {2} + \dots + a _ {2 m + 1} x _ {2 m + 1} ^ {2} = b
\]

由以下公式给出：

\(1^{\circ}\) 若 \(b = 0, \nu = q^{2m}\) ；

\(2^{\circ}\) 若 \(b \neq 0\) 且 \((-1)^{m}a_{1}a_{2} \ldots a_{2m+1}b\) 在 K 中不是平方元，则 \(\nu = q^{2m} - q^{m}\) ；

\(3^{\circ}\) 若 \(b \neq 0\) 且 \((-1)^{m}a_{1}a_{2} \ldots a_{2m+1}b\) 是 \(\mathbf{K}\) 中的平方元，则 \(\nu = q^{2m} + q^{m}\) 。（归结为 \(a\) ）的情形。）

\begin{enumerate}
\def\labelenumi{\arabic{enumi})}
\setcounter{enumi}{6}
\tightlist
\item
  设 K 为有限域，E 为 K 的 n 次代数扩张，且 \((\alpha_{1}, \ldots, \alpha_{n})\) 为 E 在 K 上的正规基；假设循环伽罗瓦群 \(\operatorname{Gal}(\operatorname{E}/\operatorname{K})\) 由循环置换 \(\sigma = (\alpha_{1}\alpha_{2} \ldots \alpha_{n})\) 生成。设 \(\tau\) 为向量 K-空间 E 的一个自同构，使得对每个 \(x \in E\) ，\(\tau(x)\) 在 E 上与 x 共轭（参见 V，第 154 页，习题 8）。
\end{enumerate}

\begin{enumerate}
\def\labelenumi{\alph{enumi})}
\item
  证明：若 \(\mathbf{K}\) 至少有 3 个元素，则 \(\tau\) 是域 E 的 K-自同构（归结到 \(\tau(\alpha_1) = \alpha_1\) 的情形，并考虑 \(\tau(\alpha_1 + c\alpha_j)\) ，其中 \(c \in \mathbf{K}\) ）。
\item
  如果 \(K = F_{2}\) ，证明对每个元素 \(\rho \in \operatorname{Gal}(E/K)\) ，关系式 \(\rho(\alpha_{r}) = \alpha_{r'}\) 与 \(\rho(\alpha_{s}) = \alpha_{s'}\) 蕴含 \(r' - r \equiv s' - s \pmod{n}\) 。由此推出，若 n \textgreater{} 5，则 \(\tau\) 是域 E 的 K-自同构。（注意：若 \(\tau(\alpha_{1}) = \alpha_{1}\) 且 \(\tau(\alpha_{1} + \alpha_{h}) = \alpha_{1} + \alpha_{k}\) ，其中 \(k \neq h\) ，则必有 \(h + k = n + 2\) ；由此推出必然有 \(\tau(\alpha_{2}) = \alpha_{n}\) 且 \(\tau(\alpha_{4}) = \alpha_{4}\) 或 \(\tau(\alpha_{4}) = \alpha_{n-2}\) ，这导致矛盾。）
\item
  对于 \(\mathbf{K} = \mathbf{F}_2\) 且 \(3 \leqslant n \leqslant 5\) ，向量 K-空间 E 的下列自同构不是 E 的域自同构：
\end{enumerate}

\[
\text { for } \quad n = 3, \tau : (\alpha_ {1}, \alpha_ {2}, \alpha_ {3}) \mapsto (\alpha_ {1}, \alpha_ {3}, \alpha_ {2})
\]

\[
\text { for } \quad n = 4, \tau : (\alpha_ {1}, \alpha_ {2}, \alpha_ {3}, \alpha_ {4}) \mapsto (\alpha_ {1}, \alpha_ {4}, \alpha_ {3}, \alpha_ {2})
\]

\[
\text { for } \quad n = 5, \tau : (\alpha_ {1}, \alpha_ {2}, \alpha_ {3}, \alpha_ {4}, \alpha_ {5}) \mapsto (\alpha_ {1}, \alpha_ {5}, \alpha_ {4}, \alpha_ {3}, \alpha_ {2})
\]

但它们满足 \(\tau(x)\) 与 \(x\) 共轭，对一切 \(x \in E\) 成立。

\begin{enumerate}
\def\labelenumi{\arabic{enumi})}
\setcounter{enumi}{7}
\tightlist
\item
  设 \(\mathbf{K}\) 是含 \(q\) 个元素的有限域，并且
\end{enumerate}

\[
\mathrm{P} (\mathbf {X}) = a _ {0} + a _ {1} \mathbf {X} + \dots + a _ {q - 2} \mathbf {X} ^ {q - 2}
\]

K{[}X{]} 中次数为 q - 2 的一个多项式。设 r 为矩阵

\[
A = \left( \begin{array}{c c c c} a _ {0} & a _ {1} & \dots & a _ {q - 2} \\ a _ {1} & a _ {2} & \dots & a _ {0} \\ \dots & \dots & \dots & \dots \\ a _ {q - 2} & a _ {0} & \dots & a _ {q - 3} \end{array} \right).
\]

证明 P 在 K 中的、异于 0 的不同根的个数是 q - 1 - r。（考虑乘积 VA，其中 V 是元素 \(x_{1}, x_{2}, \ldots, x_{q-1}\) （属于 \(K^{*}\) ）的范德蒙德矩阵。）

\begin{enumerate}
\def\labelenumi{\arabic{enumi})}
\setcounter{enumi}{8}
\tightlist
\item
  设 \(\mathbf{K}\) 是含 \(q\) 个元素的有限域，并令 \(r = q^n\) ，其中 \(n \geqslant 1\) 为整数。对每个多项式 \(\mathrm{P}(\mathbf{X}) = \sum_{j=0}^{m} a_j \mathbf{X}^j\) （属于 \(\mathbf{K}[\mathbf{X}]\) ），令
\end{enumerate}

\[
\hat {\mathrm{P}} (\mathrm{X}) = \sum_ {j} a _ {j} \mathrm{X} ^ {(r ^ {j} - 1) / (r - 1)}
\]

和

\[
\tilde {\mathrm{P}} (\mathrm{X}) = \mathrm{X} \hat {\mathrm{P}} \left(\mathrm{X} ^ {r - 1}\right) = \sum_ {j} a _ {j} \mathrm{X} ^ {r ^ {j}}.
\]

\begin{enumerate}
\def\labelenumi{\alph{enumi})}
\item
  证明：若 P、Q 是 K{[}X{]} 中的两个多项式，且 R = PQ，则我们有 \(\tilde{\mathbf{R}}(\mathbf{X}) = \tilde{\mathbf{P}}(\tilde{\mathbf{Q}}(\mathbf{X}))\) 。由此推出：要使 K{[}X{]} 中的两个多项式 F、G 满足 F 整除 G，必须而且只需 \(\hat{F}\) 整除 \(\hat{G}\) 。（要看出该条件是充分的，考虑 G 除以 F 的欧几里得除法。）
\item
  设 F 是 K{[}X{]} 中的不可约多项式，并设 G 是 K{[}X{]} 中的任意多项式。如果存在整数 \(d \geqslant 1\) 使得 \(\hat{\mathrm{F}}(\mathbf{X}^{d})\) 与 \(\hat{\mathrm{G}}(\mathbf{X}^{d})\) 有公共根，证明 F 整除 G。（若 H 是 \(\hat{\mathrm{F}}(\mathbf{X}^{d})\) 与 \(\hat{\mathrm{G}}(\mathbf{X}^{d})\) 的最大公因子，考虑多项式 \(P \in K[X]\) （H 整除 \(\hat{\mathrm{P}}(\mathbf{X}^{d})\) ）所成的集合。）
\item
  每个多项式 \(\mathbf{F} \neq 0\) （属于 \(\mathbf{K}[\mathbf{X}]\) ）都整除某个形如 \(\mathbf{X}^{\mathbf{N}} - 1\) 的多项式。证明：若 \(r - 1 = de\) 且 \(\alpha\) 是 \(\hat{\mathbf{F}}(\mathbf{X}^{d})\) 的一个根，则 \(\alpha\) 属于 \(\mathbf{K}\) 的次数为 \(n\mathbf{N}\) 的扩张，并由此推出 \(\hat{\mathbf{F}}(\mathbf{X}^d)\) 在 \(\mathbf{K}[\mathbf{X}]\) 中的每个不可约因子的次数都整除 \(n\mathbf{N}\) （利用 \(a\) ））。若进一步 \(e\) 与 \(d\mathbf{N}\) 互素，且 \(m\) 是 \(\alpha\) 在 \(\mathbf{K}\) 的次数为 \(n\) 的扩张上的次数，证明 \(m\) 整除 \(\mathbf{N}\) （利用 \((r^{\mathbf{N}} - 1)/e \equiv d\mathbf{N} (\text{mod } e)\) 这一事实）。
\end{enumerate}

\(d\) ) 采用与 \(c\) ) 相同的记号，假设 \(\mathbf{F}\) 不可约，且 \(\mathbf{N}\) 是整数 \(h\) 中最小的一个，其中 \(\mathbf{F}\) 整除 \(\mathbf{X}^h - 1\) 。若 \(e\) 与 \(d\mathbf{N}\) 互素，且 \(n\) 的每个素因子都整除 \(\mathbf{N}\) ，证明 \(\hat{\mathbf{F}}(\mathbf{X}^d)\) 的每个不可约因子的次数都是 \(n\mathbf{N}\) 。（注意，根据 \(b\) )，此时我们有 \(m = \mathbf{N}\) 。）

\begin{enumerate}
\def\labelenumi{\alph{enumi})}
\setcounter{enumi}{4}
\tightlist
\item
  特别地，由 \(d\) ) 推出：若 \(\mathrm{P}(\mathbf{X}) = \sum_{j} a_j \mathbf{X}^j\) 是
\end{enumerate}

K{[}X{]} 中的不可约多项式，且 N 是使 \(P(X)\) 整除 \(X^{h}-1\) 的最小整数 h ，则 \(\sum a_{j}X^{q^{j}-1}\) 的每个不可约因子的次数都是 N。

\begin{enumerate}
\def\labelenumi{\arabic{enumi})}
\setcounter{enumi}{9}
\tightlist
\item
  设 \(t, n, k\) 为三个 \(\geqslant 1\) 的整数，使得 \(t \leqslant n\) 且 \(n\) 整除 \(2^k - 1\) 。设 \(S\) 为满足下列条件的最小整数集合：\((1, t) \subset S \subset (1, n)\) ，且条件 \(\langle i \in S, 2i \equiv j (\text{mod } n) \text{ 且 } 1 \leqslant j \leqslant n \rangle\) 蕴含 \(j \in S\) 。令 \(m = \text{Card } S\) 。设 \(\xi\) 是一个阶为 \(n\) 的元素，含于 \(F_2^k\) ，并令 \(P(X) = \prod_{j \in S} (X - \xi^j)\) 。
\end{enumerate}

\begin{enumerate}
\def\labelenumi{\alph{enumi})}
\item
  证明 \(\mathbf{P}(\mathbf{X})\in \mathbf{F}_2[\mathbf{X}]\) 。
\item
  证明：若 \(\mathbf{R}(\mathbf{X})\in \mathbf{F}_2[\mathbf{X}]\) 的次数 \(< n\) 且是 \(\mathrm{P}(\mathbf{X})\) 的非零倍元，则其非零系数严格多于 \(t\) 个。
\item
  由此构作一个映射 \(\varphi: \mathbf{F}_2^n \to \mathbf{F}_2^m\) ，使得当 \(u, v \in \mathbf{F}_2^n\) 互异且至多有 \(t\) 个不同坐标时，有 \(\varphi(u) \equiv \varphi(v)\) 。
\item
  当 t = 6、n = 15、k = 4 时，计算 S、m 和 P。
\end{enumerate}

\begin{enumerate}
\def\labelenumi{\arabic{enumi})}
\setcounter{enumi}{10}
\tightlist
\item
  设 \(f\) 是 \(\mathbf{Z}[\mathbf{X}]\) 中的首一多项式，在 \(\mathbf{Q}[\mathbf{X}]\) 中不可约，并设 \(p\) 是一个素数，使得 \(\vec{f}\) —— \(f\) 在 \(\mathbf{F}_p[\mathbf{X}]\) 中的典范像 —— 是一个无重根的多项式。若 \(\Gamma\) （相应地 \(\vec{\Gamma}\) ）是多项式 \(f\) （相应地 \(\vec{f}\) ）在 \(\mathbf{Q}\) （相应地 \(\mathbf{F}_p\) ）上的伽罗瓦群，试给出从 \(\Gamma\) 的一个子群到 \(\vec{\Gamma}\) 之上的一个同构。由此推出：若 \(\vec{f} = g_1g_2\ldots g_r\) ，其中各 \(g_j\) 在 \(\mathbf{F}_p[\mathbf{X}]\) 中不可约，且 \(g_j\) 的次数为 \(n_j\) （从而 \(f\) 的次数为 \(n = n_1 + n_2 + \dots + n_r\) ），则群 \(\Gamma\) 作为 \(\mathfrak{S}_n\) 的子群，包含一个由支集两两不相交的循环组成的乘积 \(\sigma_1\sigma_2\ldots\sigma_r\) ，其中 \(\sigma_j\) 的阶为 \(n_j\) ，\(1\leqslant j\leqslant r\) （参见 I，第 62 页）。
\end{enumerate}

推广到 f 不是首一多项式、但其首项系数不被 p 整除的情形。

\begin{enumerate}
\def\labelenumi{\arabic{enumi})}
\setcounter{enumi}{11}
\item
  设 \(f\) 是 \(\mathbf{Z}[\mathbf{X}]\) 中次数为 \(n\) 的多项式，并设 \(p_1, p_2, p_3\) 是三个不整除 \(f\) 的首项系数的互异素数；设 \(f_1, f_2, f_3\) 分别是 \(f\) 在 \(\mathbf{F}_{p_1}[\mathbf{X}], \mathbf{F}_{p_2}[\mathbf{X}], \mathbf{F}_{p_3}[\mathbf{X}]\) 中的典范像。假设 \(f_1\) 是一个线性因子与一个次数为 \(n - 1\) 的不可约因子之积，\(f_2\) 是一个二次不可约因子与若干奇数次不可约因子之积，最后 \(f_3\) 不可约。证明：在这些条件下，\(f\) 在 \(\mathbf{Q}[\mathbf{X}]\) 中不可约，且多项式 \(f\) 在 \(\mathbf{Q}\) 上的伽罗瓦群同构于对称群 \(\mathfrak{S}_n\) （«戴德金准则»；同时利用习题 11 以及 I，第 63 页，命题 9）。
\item
  设 \(p_1, p_2, \ldots, p_m\) 为互不相同的奇素数，且 \(P = p_1p_2 \ldots p_m\) （ \(m \geqslant 3\) ）。
\end{enumerate}

\begin{enumerate}
\def\labelenumi{\alph{enumi})}
\tightlist
\item
  给定整数 \(n \geqslant 1\) ，设 \(E\) 为 \(\mathbf{Z}[\mathbf{X}]\) 中次数 \(\leqslant n\) 、系数落在区间 \([0, P[\) （属于 \(\mathbf{N}\) ）内的多项式组成的集合；于是我们有 \(\text{Card}(E) = P^{n+1}\) 。设 \(k\) 是满足 \(0 < k < 1/3n\) 的数；证明：在多项式 \(f \in E\) 之中，至少有 \(k^m\mathbf{P}^{n + 1}\) 个使得对每个指标 \(j\) （满足 \(1 \leqslant j \leqslant m\) ），典范像 \(f_j\) —— \(f\) 在 \(\mathbf{F}_{p_j}[\mathbf{X}]\) 中的 —— 的次数为 \(n\) 且不可约（利用习题 1，b））。再假设 \(k < \frac{1}{18(n - 2)}\) 且 \(n > 2\) ；以同样的方式论证并利用戴德金准则（习题 12）
\end{enumerate}

证明存在一个由多项式 \(f \in \mathbf{E}\) 组成的集合，其基数至少等于 \((1 - 3(1 - k)^m) \mathbf{P}^{n+1}\) ，这些多项式次数为 \(n\) 、不可约，且其在 \(\mathbf{Q}\) 上的伽罗瓦群同构于 \(\mathfrak{S}_n\) 。

\begin{enumerate}
\def\labelenumi{\alph{enumi})}
\setcounter{enumi}{1}
\tightlist
\item
  对每个整数 N，设 \(L_{n,N}\) 为 Z{[}X{]} 中次数 \(\leqslant n\) 、系数属于区间 \((-N, N)\) 的多项式组成的集合；于是我们有 \(\text{Card}(L_{n,N}) = (2N + 1)^{n+1}\) 。固定整数 n，证明：对每个 \(\varepsilon\) （满足 \(0 < \varepsilon < 1\) ），存在整数 \(N_0\) ，使得当 \(N \geqslant N_0\) 时，多项式 \(f \in L_{n,N}\) 中至少有 \((1 - \varepsilon)(2N + 1)^{n+1}\) 个是次数为 n 的不可约多项式，且其在 Q 上的伽罗瓦群同构于 \(S_n\) （利用 a））。
\end{enumerate}

* 14) 设 K 是一个有限域，交换与否均可；设 Z 是它的中心，q 是 Z 的元素个数，n 是秩 {[}K:Z{]} 。

\begin{enumerate}
\def\labelenumi{\alph{enumi})}
\item
  若 E 是 K 的包含 Z 的子域，证明 \([E:Z]\) 整除 n。
\item
  设 \(x \in K\) 是不属于 Z 的元素；证明 x 的不同共轭元 \(yxy^{-1}\) 在 \(K^{*}\) 中的个数具有形式 \((q^{n}-1)/(q^{d}-1)\) ，其中 d 整除 n 且异于 n（在 K 中考虑与 x 可交换的元素组成的集合，并利用 a））。
\item
  由 \(b\) ) 推出 \(q - 1\) 被整数 \(\Phi_n(q)\) 整除（把群 \(\mathbf{K}^*\) 划分成共轭类，并利用 \(V\) ，第 82 页的关系式 (6)）。
\item
  证明：若 \(n > 1\) ，则 \(\Phi_n(q) > (q - 1)^{\varphi(n)}\) （把 \(\Phi_n(X)\) 在复数域 \(C\) 中分解）。由此推断必有 \(K = Z\) ，换言之，\(K\) 必然是交换的（韦德伯恩定理）。*
\end{enumerate}

\BourbakiExerciseGroup

\begin{enumerate}
\def\labelenumi{\arabic{enumi})}
\tightlist
\item
  设 K 是特征为 p \textgreater{} 0 的域，E 是 K 的一个扩张；E 是一个高度 \(\leqslant 1\) 的、\(\mathrm{K}(\mathrm{E}^{p})\) 的 p-根式扩张。E 在 \(\mathrm{K}(\mathrm{E}^{p})\) 上的 p-基的基数（参见 V，第 103 页，定理 2）称为 E 在 K 上的不完全度。若 \(\mathrm{K}(\mathrm{E}^{p}) = \mathrm{E}\) ，则称 E 在 K 上相对完全；若 E 是完全域，则它在其每个子域上都相对完全。K 的（绝对）不完全度是 K 在其素子域 \(F_{p}\) 上的不完全度，亦即 K 的一个（绝对）p-基的基数。
\end{enumerate}

\begin{enumerate}
\def\labelenumi{\alph{enumi})}
\item
  若 B 是 E 的一个 \(p\) -基（在 \(\mathbf{K}(\mathbf{E}^p)\) 上），证明对每个整数 \(k > 0\) ，我们有 \(\mathbf{E} = \mathbf{K}(\mathbf{E}^{p^k})(\mathbf{B})\) 。
\item
  假设 \(\mathbf{E} \subset \mathbf{K}^{p^{-n}}\) 对某个整数 \(n > 0\) 成立。证明：要使次数 \([\mathbf{E}:\mathbf{K}]\) 有限，必须而且只需不完全度 \(m_0\) —— \(\mathbf{E}\) 在 \(\mathbf{K}\) 上的 —— 为有限；此时 \(m_0\) 是集合 \(\mathbf{S}\) （满足 \(\mathbf{E} = \mathbf{K}(\mathbf{S})\) ）的最小基数。证明：若 \(m_k\) 是 \(\mathbf{K}(\mathrm{E}^{p^k})\) 在 \(\mathbf{K}\) 上的不完全度，则 \(m_{k + 1} \leqslant m_k\) 对所有 \(k\) 成立，且若 \(f = \sum_{k} m_k\) ，则 \([\mathbf{E}:\mathbf{K}] = p^f\) 。
\end{enumerate}

\begin{enumerate}
\def\labelenumi{\arabic{enumi})}
\setcounter{enumi}{1}
\tightlist
\item
  \begin{enumerate}
  \def\labelenumii{\alph{enumii})}
  \tightlist
  \item
    设 E 是特征为 p \textgreater{} 0 的域 K 的扩张，F 是 E 的扩张。证明：若 B 是 E 在 \(\mathrm{K}(\mathrm{E}^{p})\) 上的 p-基，C 是 F 在 \(\mathrm{E}(\mathrm{F}^{p})\) 上的 p-基，则存在 F 在 \(\mathrm{K}(\mathrm{F}^{p})\) 上的一个包含于 \(B \cup C\) 的 p-基。
  \end{enumerate}
\end{enumerate}

\begin{enumerate}
\def\labelenumi{\alph{enumi})}
\setcounter{enumi}{1}
\tightlist
\item
  假设 B 有限，且 F 是 K 的有限次代数扩张。证明 \([E:K(E^{p})]\geqslant[F:K(F^{p})]\) 。若进一步 F 在 E 上可分，证明 \([E:K(E^{p})]=[F:K(F^{p})]\) 。
\end{enumerate}

\begin{enumerate}
\def\labelenumi{\arabic{enumi})}
\setcounter{enumi}{2}
\item
  设 K 是一个非完全域，E 是 K 的有限次代数扩张。要使存在元素 \(x \in E\) 使得 \(\mathrm{E} = \mathrm{K}(x)\) ，必须而且只需 E 在 K 上的不完全度（习题 1）为 0 或 1。（要看出该条件是充分的，注意：若 \(E_{s}\) 是 K 在 E 中的相对可分闭包，我们有 \(\mathrm{E} = \mathrm{E}_{s}(\alpha)\) 且 \(\mathrm{E}_{s} = \mathrm{K}(\beta)\) ；若 \(\alpha \notin E_{s}\) ，考虑 \(\alpha + c\beta\) 在 K 上的共轭元，其中 \(c \in K\) 。）
\item
  设 K 是特征为 p \textgreater{} 0 的域，E 是 K 的有限次代数扩张。若 r \textgreater{} 0 是 E 在 K 上的不完全度（习题 1），证明 r 是使 \(\mathrm{E} = \mathrm{K}(\mathrm{S})\) 的集合 S 的最小基数。（要看出 E 可由它的 r 个元素生成，注意：若 \(E_{s}\) 是 K 在 E 中的相对可分闭包，则存在 E 的 r 个元素 \(a_{i}\) （ \(1 \leqslant i \leqslant r\) ）使得 \(\mathrm{E} = \mathrm{E}_{s}(a_{1}, \ldots, a_{r})\) ，并利用习题 3。）
\item
  设 K 是一个域，E 是 K 的有限次代数扩张。
\end{enumerate}

\begin{enumerate}
\def\labelenumi{\alph{enumi})}
\item
  假设 K 的特征为 p \textgreater{} 0。证明：若 E 在 K 上的不完全度（习题 1）\textgreater{} 1，则存在无穷多个不同的域 F，使得 \(K \subset F \subset E\) （归结为 \(\mathrm{K}(\mathrm{E}^{p}) = \mathrm{K}\) 的情形，并考虑域 \(\mathrm{K}(a + \lambda b)\) ，其中集合 \(\{a, b\}\) 在 K 上是 p-自由的，且 \(\lambda \in K\) ）。
\item
  由此得出结论：要使满足 \(K \subset F \subset E\) 的域 F 只有有限多个，必须而且只需 E 在 K 上的不完全度为 \(\leqslant 1\) 。
\end{enumerate}

\begin{enumerate}
\def\labelenumi{\arabic{enumi})}
\setcounter{enumi}{5}
\tightlist
\item
  设 K 是特征为 p \textgreater{} 0 的域，E 是 K 的有限次代数扩张，N 是由 E 生成的 K 的拟伽罗瓦扩张。
\end{enumerate}

\begin{enumerate}
\def\labelenumi{\alph{enumi})}
\item
  设 \(E_{r}\) 与 \(E_{s}\) （相应地 \(N_{r}\) 与 \(N_{s}\) ）是包含在 E（相应地 N）中的 K 的最大 p-根式扩张与最大可分扩张。证明：要使 \(\mathrm{E} = \mathrm{E}_{r}(\mathrm{E}_{s})\) （此时 E 同构于 \(E_{s} \otimes_{K} E_{r}\) ），必须而且只需 \(E_{r} = N_{r}\) 。
\item
  若 E 在 K 上的不完全度为 1，则 \(\mathrm{E} = \mathrm{E}_{r}(\mathrm{E}_{s})\) 的充分必要条件是 N 在 K 上的不完全度为 1。由此推出：若 N 是 K 的拟伽罗瓦扩张，且其在 K 上的不完全度等于 1，则对 N 的每个子扩张 E 都有 \(\mathrm{E} = \mathrm{E}_{r}(\mathrm{E}_{s})\) （参见习题 2，b））。
\item
  反之，若 N 是 K 的有限次拟伽罗瓦扩张，且其在 K 上的不完全度 \textgreater1，证明：若 \(N_{r} \neq N\) ，则存在 \(t \in N\) 使得对 K 的扩张 \(\mathrm{E} = \mathrm{K}(t)\) 我们有 \(\mathrm{E} \neq \mathrm{E}_{r}(\mathrm{E}_{s})\) （若 a、b 是 N 在 \(\mathrm{K}(\mathrm{N}^{p})\) 上的一个 p-基中的两个元素，且 \(x \in N\) 在 K 上可分并满足 \(x \notin K\) ，取 \(t = a + bx\) ）。
\end{enumerate}

\BourbakiExerciseGroup

\begin{enumerate}
\def\labelenumi{\arabic{enumi})}
\item
  设 E 是域 K 的一个扩张，B 是 E 在 K 上的一个超越基。证明：若集合 K、B 中有一个是无限的，则 E 与 \(K \times B\) 等势，否则 E 是可数的。* 特别地，由此推出实数域 R 在有理数域 Q 上的每个超越基都具有连续统的势。*
\item
  设 K 为 \(\mathbf{Q}(\mathbf{X})\) ，即有理数域 Q 上一个未定元的有理分式域。证明：在环 K{[}Y{]} 中，多项式 \(Y^{2} + X^{2} + 1\) 不可约；又若 E 是由该多项式的一个根（在
\end{enumerate}

K），则 Q 在 E 中相对代数封闭，但 E 不是 Q 的纯扩张。（要看出 E 中没有在 Q 上代数的元素 \(a \notin Q\) ，注意这样一个元素的存在将蕴含关系 \(\mathrm{E} = \mathbf{Q}(a)(\mathbf{X})\) ，并注意 \(-(\mathbf{X}^{2} + 1)\) 在 \(\Omega(\mathbf{X})\) 中不是平方元，其中 \(\Omega\) 是 Q 的一个代数闭包。）证明：若 i 是 \(X^{2} + 1\) 的一个根，则 \(\mathrm{E}(i)\) 是 \(\mathbf{Q}(i)\) 的纯超越扩张。

* 3) 设 K 为域 C(X)，即（代数闭的）复数域 C 上一个未定元的有理分式域。证明：在环 K{[}Y{]} 中，多项式 \(Y^{3} + X^{3} + 1\) 不可约；设 E 为由该多项式的一个根（在 K 的一个代数闭包中）生成的 K 的扩张。证明 E 不是 C 的纯扩张，尽管 C 是代数闭的。（证明：关系 \(u^{3} + v^{3} + w^{3} = 0\) 不可能存在于 C{[}X{]} 的三个两两互素且并非全为常数的多项式 u、v、w 之间。用反证法论证：若 r 是 u、v、w 的次数中的最大者，且例如 \(\deg(w) = r\) ，则由关系式

\[
w ^ {3} = - (u + v) (u + j v) (u + j ^ {2} v)
\]

其中 \(j\) 是本原三次单位根，由该关系式推出存在三个多项式 \(u_1\) 、\(v_1\) 、\(w_1\) （在 \(\mathbf{C}[\mathbf{X}]\) 中），它们两两互素、并非全为常数、次数 \(< r\) ，且满足 \(u_1^3 + v_1^3 + w_1^3 = 0\) 。）*

\begin{enumerate}
\def\labelenumi{\arabic{enumi})}
\setcounter{enumi}{3}
\tightlist
\item
  设 \(\Omega\) 是域 \(\mathbf{K}\) 的一个代数闭扩张，且在 \(\mathbf{K}\) 上具有无限超越次数。
\end{enumerate}

\begin{enumerate}
\def\labelenumi{\alph{enumi})}
\tightlist
\item
  证明：存在无穷多个 \(\Omega\) 的 K-自同态，其像是 \(\Omega\) 的异于 \(\Omega\) 的子域，并且关于这些自同态，\(\Omega\) 可以具有等于直至 tr. \(\deg_{\mathbf{K}}\Omega\) 的任意基数的超越次数。* 特别地，存在无穷多个不同的复数域 C 到 C 的异于 C 的子域上的 Q-同构。* b) 证明：对 \(\Omega\) 的每个满足 tr. \(\deg_{\mathbf{K}}E < \text{tr} \cdot \deg_{\mathbf{K}}\Omega\) 的子扩张 E，E 到 \(\Omega\) 的一个子域上的每个 K-同构都可以延拓为 \(\Omega\) 的一个 K-自同构。给出这样的例子：其中 tr. \(\deg_{\mathbf{K}}E = \text{tr} \cdot \deg_{\mathbf{K}}\Omega\) ，且存在 E 到 \(\Omega\) 的一个子域上的一个 K-同构，它不能延拓为 E 的任一扩张 F（包含于 \(\Omega\) 且异于 E）到 \(\Omega\) 的一个子域上的 K-同构。
\end{enumerate}

\begin{enumerate}
\def\labelenumi{\arabic{enumi})}
\setcounter{enumi}{4}
\item
  设 \(\Omega\) 是域 \(\mathbf{K}\) 的一个代数闭扩张。证明：对每个子扩张 \(\mathbf{E}\) —— \(\Omega\) 的、在 \(\mathbf{K}\) 上为超越的 —— 存在无穷多个 \(\mathbf{K}\) -同构，把 \(\mathbf{E}\) 映到 \(\Omega\) 的一个子域上（若 \(x \in \mathbf{E}\) 在 \(\mathbf{K}\) 上是超越的，考虑子扩张 \(\mathbf{F}\) （含于 \(\mathbf{E}\) ），使得 \(x\) 在 \(\mathbf{F}\) 上超越且 \(\mathbf{E}\) 在 \(\mathbf{F}(x)\) 上代数，并证明存在无穷多个 \(\mathbf{F}\) -同构，从 \(\mathbf{E}\) 到 \(\Omega\) 的一个子域上）。
\item
  设 \(\Omega\) 是域 \(K\) 的一个代数闭扩张，\(N\) 是 \(\Omega\) 的一个子扩张。证明：若 \(N\) 在 \(K\) 上是拟伽罗瓦的，且 \(E\) 与 \(E'\) 是 \(K\) 的两个包含在 \(\Omega\) 中的共轭扩张，则 \(N(E)\) 与 \(N(E')\) 在 \(K\) 上共轭。反之，若 \(N\) 具有此性质，且 \(\operatorname{tr} \cdot \deg_{K} N\) 有限并 \(< \operatorname{tr} \cdot \deg_{K} \Omega\) ，则它必然在 \(K\) 上是代数的且是拟伽罗瓦的。
\item
  设 E 是域 K 的一个扩张，\(\operatorname{Aut}_{\mathbf{K}}(\mathbf{E})\) 是 E 的所有 K-自同构组成的群。a) 对 E 在 K 上的每个有限生成子扩张 L，\(\operatorname{Aut}_{\mathbf{L}}(\mathbf{E})\) 是 \(\operatorname{Aut}_{\mathbf{K}}(\mathbf{E})\) 的一个子群。证明：当 L 遍历 E 的在 K 上有限生成的子扩张所成的集合时，诸群 \(\operatorname{Aut}_{\mathbf{L}}(\mathbf{E})\) 构成 \(\mathrm{Aut}_{\mathbf{K}}(\mathbf{E})\) 的单位元的一个基本邻域系，对应于一个与 \(\mathrm{Aut}_{\mathbf{K}}(\mathbf{E})\) 的群结构相容的拓扑，且对该拓扑 \(\mathrm{Aut}_{\mathbf{K}}(\mathbf{E})\) 是一个分离且完全不连通的群。
\end{enumerate}

\begin{enumerate}
\def\labelenumi{\alph{enumi})}
\setcounter{enumi}{1}
\tightlist
\item
  当 E 赋有离散拓扑时，\(\operatorname{Aut}_{\mathbf{K}}(\mathbf{E})\) 的拓扑是使映射 \((u,x)\to u(x)\) —— 从 \(\operatorname{Aut}_{\mathbf{K}}(\mathbf{E})\times\mathbf{E}\) 到 E 的映射 —— 连续的最粗拓扑。
\end{enumerate}

* c) 群 \(\operatorname{Aut}_{\mathbf{Q}}(\mathbf{R})\) 约化为中性元素（V，第 153 页，§9 习题 2）。* \(d)\) 设 \(\mathbf{K}_0 \supseteq \mathbf{K}\) 是 E 的由在 E 的所有 K-自同构下不变的元素组成的子扩张。证明：要使 E 在 \(\mathbf{K}_0\) 上是代数的，必须而且只需 \(\operatorname{Aut}_{\mathbf{K}}(\mathbf{E})\) 是紧致的。

\begin{enumerate}
\def\labelenumi{\alph{enumi})}
\setcounter{enumi}{4}
\tightlist
\item
  假设 K 是无限的。若 \(E = K(X)\) 是 K 上一个未定元的有理分式域，证明 \(\operatorname{Aut}_{K}(E)\) 是一个无限离散群（V，第 148 页，习题 11）。若 K 的特征为 0，则 \(\Gamma\) —— \(\operatorname{Aut}_{K}(E)\) 中由自同构 \(\sigma : X \mapsto X + 1\) 生成的子群 —— 在 \(\operatorname{Aut}_{K}(E)\) 中是闭的且异于后者，但具有相同的不变域。
\end{enumerate}

\begin{enumerate}
\def\labelenumi{\arabic{enumi})}
\setcounter{enumi}{7}
\tightlist
\item
  设 \(E\) 是域 \(K\) 的一个扩张。
\end{enumerate}

\begin{enumerate}
\def\labelenumi{\alph{enumi})}
\item
  设 \(x\) 是 \(E - K\) 的一个元素。证明：存在一个纯扩张 \(L \subset E\) （ \(K\) 的），使得 \(E\) 在 \(L\) 上是代数的，且 \(x \notin L\) （考虑子扩张 \(L\) —— 属于 \(E\) 、在 \(K\) 上纯超越且 \(x \notin L\) 的 —— 所成的集合）。
\item
  设 \(x \in E\) 是 \(K\) 上的一个超越元，\(p\) 是一个整数 \(\geqslant 2\) 。证明：存在一个纯超越扩张 \(L \subset E\) （ \(K\) 的），使得 \(E\) 在 \(L\) 上是代数的，\(x \notin L\) 且 \(x^p \in L\) （同样的方法）。
\item
  设 F 是 E 的一个子扩张，\(x \in E\) 是 F 上的一个代数元，\(P \in F[X]\) 是它的极小多项式。证明存在 K 的一个纯扩张 \(L \subset E\) ，使得 E 在 L 上是代数的，且 P 在域 \(F(L)\) 上不可约。
\end{enumerate}

\begin{enumerate}
\def\labelenumi{\arabic{enumi})}
\setcounter{enumi}{8}
\tightlist
\item
  域 K 的一个扩张 E 称为 K 的戴德金扩张，如果对 E 的每个子扩张 F，F 都与 E 中在群 \(\mathrm{Aut}_{\mathrm{F}}(\mathrm{E})\) —— E 的所有 F-自同构组成的群 —— 作用下不变的元素所构成的子域相同。
\end{enumerate}

\begin{enumerate}
\def\labelenumi{\alph{enumi})}
\item
  要使 \(\mathbf{E}\) —— \(\mathbf{K}\) 的代数扩张 —— 在 \(\mathbf{K}\) 上是戴德金的，必须而且只需 \(\mathbf{E}\) 是 \(\mathbf{K}\) 的一个伽罗瓦扩张。
\item
  假设 K 的特征为 p \textgreater{} 0。证明 K 的戴德金扩张必然在 K 上是代数的（从而是伽罗瓦的）。（利用习题 8，b）。）
\item
  如果E是K的戴德金扩张，且L是E的一个纯超越子扩张，不同于K，并且E是L的代数扩张，证明E在L上的次数是无限的。
\end{enumerate}

\begin{enumerate}
\def\labelenumi{\arabic{enumi})}
\setcounter{enumi}{9}
\item
  设 E 是域 K 的一个扩张，使得映射 \(F \mapsto \operatorname{Aut}_{F}(E)\) 是从 E 的子扩张所成的集合到 \(\operatorname{Aut}_{K}(E)\) 的子群所成的集合上的一个双射。证明此时 E 是 K 的一个有限次伽罗瓦扩张。（利用习题 9，c)，先证明 E 在 K 上必然是代数的，再利用习题 9，a)。）
\item
  设 \(\mathbf{E}\) 是域 \(\mathbf{K}\) 的一个扩张，\(\Omega\) 是 \(\mathbf{E}\) 的一个代数闭扩张。证明下列条件是等价的：
\end{enumerate}

α) E是K的p-根式扩张；

\(\beta)\) 对每个扩张 \(F\) （在 \(K\) 之上），\(E \otimes_{K} F\) 只有唯一的素理想；

γ) 对于K的每个扩张F，E与F的任意两个复合扩张都是同构的；

δ) \(E \otimes_{K} \Omega\) 只有一个素理想；

ε) E 和 Ω 的任意两个复合扩张都是同构的。

\begin{enumerate}
\def\labelenumi{\arabic{enumi})}
\setcounter{enumi}{11}
\tightlist
\item
  \begin{enumerate}
  \def\labelenumii{\alph{enumii})}
  \tightlist
  \item
    设 K 是一个域，\(\Omega\) 是 K 的一个代数闭扩张，\(E \subset \Omega\) 与 F 是 K 的两个扩张；假设 \(\operatorname{tr} \cdot \deg_E \Omega \geqslant \operatorname{tr} \cdot \deg_K F\) 。证明 E 与 F 的每个复合扩张都同构于形如 \((L_w, 1, w)\) 的复合，其中 \(w\) 是 F 到 \(\Omega\) 的一个子域上的一个 K-同构，而 \(L_w\) 是 \(\Omega\) 中由 E 与 \(w(F)\) 生成的子域。
  \end{enumerate}
\end{enumerate}

\begin{enumerate}
\def\labelenumi{\alph{enumi})}
\setcounter{enumi}{1}
\tightlist
\item
  由 a) 推出：若 F 在 K 上是 n 次有限代数扩张，则 E 与 F 的复合扩张中两两不同构的至多有 n 个。
\end{enumerate}

\begin{enumerate}
\def\labelenumi{\arabic{enumi})}
\setcounter{enumi}{12}
\item
  在 \(\Omega\) —— \(\mathbf{Q}(\mathbf{X})\) 的一个代数闭包 —— 中，考虑 \(\mathrm{E} = \mathbf{Q}(\mathbf{X})\) 与 \(\mathrm{F} = \mathbf{Q}(\mathbf{X} + i)\) （其中 \(i^2 = -1\) ）这两个域 \(\mathbf{Q}\) 的纯超越扩张。证明 \(\mathrm{E} \cap \mathrm{F} = \mathbf{Q}\) ，但 \(\mathrm{E}\) 与 \(\mathrm{F}\) 在 \(\mathbf{Q}\) 上并非代数不相交。
\item
  设 E、F、G 是域 K 的包含于 K 的一个扩张 \(\Omega\) 中的三个扩张，且满足 \(F \subset G\) 。证明：要使 E 与 G 在 K 上代数不相交，必须而且只需 E 与 F 在 K 上代数不相交，并且 \(\mathrm{E}(\mathrm{F})\) 与 G 在 F 上代数不相交。
\item
  \begin{enumerate}
  \def\labelenumii{\alph{enumii})}
  \tightlist
  \item
    设 K 是一个域，L 是 K 的一个子域，使得 K 是一个有限生成的 L-代数，换言之，\(K = L[a_{1}, a_{2}, \ldots, a_{n}]\) ，其中元素属于 K。证明这些 \(a_{j}\) 在 L 上是代数的。（用反证法论证：若 \(a_{1}, \ldots, a_{m}\) （ \(m \geqslant 1\) ）构成 \((a_{j})_{1 \leqslant j \leqslant n}\) 的一个极大代数无关子族，注意在环 \(A = L[a_{1}, \ldots, a_{m}]\) 中所有非零理想的交集为零，并利用 V，第 147 页，习题 5。）
  \end{enumerate}
\end{enumerate}

\begin{enumerate}
\def\labelenumi{\alph{enumi})}
\setcounter{enumi}{1}
\tightlist
\item
  由 a) 推出：在多项式代数 \(K[X_{1}, \ldots, X_{n}]\) 中，每个极大理想在 K 上的余维数都是有限的。由此得出结论：若 A 是 K 上的一个交换有限生成代数，则 A 的包含 K 的每个子域在 K 上的次数都是有限的。
\end{enumerate}

\begin{enumerate}
\def\labelenumi{\arabic{enumi})}
\setcounter{enumi}{15}
\tightlist
\item
  \begin{enumerate}
  \def\labelenumii{\alph{enumii})}
  \tightlist
  \item
    设 K 是一个代数闭域，\(\Omega\) 是 K 的一个代数闭扩张，E 是 \(\Omega\) 的一个子扩张，\(x\) 是 \(\Omega\) 的一个在 E 上超越的元素。设 \(\overline{\mathbf{K}(x)}\) 是 \(\mathbf{K}(x)\) 在 \(\Omega\) 中的相对代数闭包，并令 \(\mathbf{M} = \operatorname{E}(\overline{\mathbf{K}(x)})\) ；证明 M 是 \(\operatorname{E}(x)\) 的一个无限次代数扩张。（注意，对每个整数 \(m > 1\) ，多项式 \(\mathbf{X}^m - x\) 在 \(\operatorname{E}(x)[\mathbf{X}]\) 中不可约（V，第 91 页，注记）。
  \end{enumerate}
\end{enumerate}

\begin{enumerate}
\def\labelenumi{\alph{enumi})}
\setcounter{enumi}{1}
\tightlist
\item
  设 E 是一个代数闭域，F 是 E 的一个有限生成扩张。证明 F 的每个代数闭子域必然包含在 E 中。（否则将存在一个代数闭的子域 \(L \subset F\) ，以及一个在 E 上超越的元素 \(x \in L\) 。应用 a)，取 K 为 E 的素子域在 E 中的相对代数闭包，并注意 M 应在 E 上有限生成。）
\end{enumerate}

\begin{enumerate}
\def\labelenumi{\arabic{enumi})}
\setcounter{enumi}{16}
\item
  设 K 是一个域，\(\Omega\) 是 K 的一个代数闭扩张，x 是 \(\Omega\) 的一个在 K 上超越的元素。证明：对每个素数 q，存在 \(\Omega\) 的一个子扩张 M，使得 \(\mathbf{M}(x)\) 是 M 的 q 次伽罗瓦扩张。（若 q 不是 K 的特征，考虑元素 \(y = x^{q}\) ，在相反的情形考虑 \(y = x^{q} - x\) ；再考虑 \(\Omega\) 的一个子扩张 L，它由 \(\Omega\) 在 K 上的一个包含 y 的超越基生成，且若 \(y = x^{q}\) 则还加上 q 次单位根；利用 V，第 163 页，习题 12。）
\item
  设 K 是一个域，\(\mathbf{K}((\mathbf{X}))\) 是 K 上一个未定元的形式幂级数域（IV，第 38 页）。证明 \(\operatorname{tr}.\deg_{\mathbf{K}}\mathbf{K}((\mathbf{X}))=\operatorname{Card}(\mathbf{K}^{\mathbf{N}})\) 。区分两种情形：
\end{enumerate}

\begin{enumerate}
\def\labelenumi{\alph{enumi})}
\item
  若 \(\operatorname{Card}(\mathbf{K}) < \operatorname{Card}(\mathbf{K}^{\mathbf{N}})\) ，注意 \(\operatorname{Card}(\mathbf{K}((\mathbf{X}))) = \operatorname{Card}(\mathbf{K}^{\mathbf{N}})\) 。
\item
  若 \(\operatorname{Card}(K) = \operatorname{Card}(K^{N})\) ，设 P 是 K 的素子域，S 是 \(\operatorname{K}((X))\) 中一个在 K 上代数无关的无限元素集，T 是属于 S 的所有形式幂级数的系数所成的集合。设 L 是域 K(S) 在
\end{enumerate}

K((X)) 中的相对代数闭包，并设 u 是 L 的一个元素；存在一个方程 \(g(s_{1}, \ldots, s_{m}, u) = 0\) ，其中 \(g \neq 0\) 是 \(K[X_{1}, \ldots, X_{m}, X_{m+1}]\) 中的一个多项式，而 \(s_{1}, \ldots, s_{m}\) 是 S 的元素。设 A 是多项式 g 的系数所成的集合，\(C(u)\) 是形式幂级数 u 的系数所成的集合；证明域 \(P(T)(A \cup C(u))\) 在 \(P(T \cup A)\) 上是代数的（用 \(\Omega\) 表示 K 的一个代数闭包，证明否则将存在无穷多个形式幂级数 \(v \in \Omega((X))\) 满足方程 g（ \(s_{1}, \ldots, s_{m}, v) = 0\) ）。利用习题 1，证明若 \(\text{Card}(S) < \text{Card}(K)\) ，则 K 在 \(P(T)\) 上的超越次数是无限的，并由此推出在这种情形 L 异于 \(K((X))\) （若 \((t_{n})_{n \geqslant 0}\) 是 K 中在 \(P(T)\) 上代数无关的一个无穷序列，考虑形式幂级数 \(\sum_{n=0}^{\infty} t_{n} X^{n}\) ）。

\begin{enumerate}
\def\labelenumi{\arabic{enumi})}
\setcounter{enumi}{18}
\tightlist
\item
  设 \(\mathbf{E}\) 与 \(\mathbf{F}\) 是域 \(\mathbf{K}\) 的两个在 \(\mathbf{K}\) 上线性不相交的超越扩张。证明 \(\mathbf{K}(\mathbf{E} \cup \mathbf{F})\) 异于环 \(\mathbf{C}\) （同构于 \(\mathbf{E} \otimes_{\mathbf{K}} \mathbf{F}\) ，由 \(\mathbf{E} \cup \mathbf{F}\) 生成的）。（归结为 \(\mathbf{E}\) 与 \(\mathbf{F}\) 在 \(\mathbf{K}\) 上的超越次数都等于 1 的情形；若 \(x \in \mathbf{E}\) 与 \(y \in \mathbf{F}\) 在 \(\mathbf{K}\) 上是超越的，证明我们不可能有 \((x + y)^{-1} \in \mathbf{C}\) ；用反证法论证，证明在相反的情形，若 \(p\) 是 \(\mathbf{K}\) 的特征指数，则将存在整数 \(r \geqslant 0\) 使得 \((x + y)^{-p^r}\) 属于 \(\mathbf{C}\) 的由 \(\mathbf{K}(x) \cup \mathbf{K}(y)\) 生成的子环。）
\end{enumerate}

* 20) 设 \(n\) 是一个整数，\(K\) 是一个特征 \(p\) 与 \(n\) 互素、且具有本原 \(n\) 次单位根的域，\(\mu_n \subset K^*\) 是 \(n\) 次单位根所成的子群，\(G\) 是一个被 \(n\) 零化的有限交换群，而 \(A = K^{(G)}\) 是 \(G\) 的群代数（III，第 446 页）。

\begin{enumerate}
\def\labelenumi{\alph{enumi})}
\item
  证明 A 是一个可对角化代数（V，第 28 页）。（将 G 分解为循环群的直和（VII，第 22 页）。）由此推出：每个 A-模都是一些模的直和，这些模是维数为 1 的向量 K-空间。
\item
  设 V 是 K 上有限维的 A-模，S 是向量 K-空间 V 的对称代数，F 是 S 的分式域。群 G 作用在 V 上，从而作用在 S 与 F 上。假设 G 在 V 上的作用是忠实的。证明 G 在 F 上的作用是忠实的。计算 \([F : F^{G}]\) 。
\item
  设 \((x_{1},\ldots ,x_{n})\) 是向量 K-空间 V 的一组由 G 的作用的特征向量组成的基。设 \(\Gamma\) 是 \(\mathbf{F}^*\) 的由 \(x_{1},\ldots ,x_{n}\) 生成的子群。证明 \(\Gamma\) 是一个自由交换群，且 F 的由 \(\Gamma\) 生成的子 K-代数 B 可与 \(\mathbf{K}^{(\Gamma)}\) 等同，并且在 G 的作用下是稳定的。
\item
  对所有 \(x \in \Gamma\) ，映射 \(\gamma \mapsto \frac{\gamma x}{x}\) 是从 \(G\) 到 \(\mu_n\) 中的一个同态，由此给出一个同态 \(\theta: \Gamma \to \operatorname{Hom}(G, \mu_n)\) 。证明 \(\theta\) 是满射的。令 \(\Gamma_1 = \operatorname{Ker}(\theta)\) 。利用初等因子定理（VII，第 24 页），证明 \(B\) 是一个秩为 Card(G) 的自由 \(K^{(\Gamma_1)}\) -模。
\item
  通过过渡到分式域，由 \(d\) 与 \(a\) 推出 \(\mathbf{F}^{\mathrm{G}}\) 是 \(\mathbf{K}^{(\Gamma_1)}\) 的分式域。由此推出 \(\mathbf{F}^{\mathrm{G}}\) 是 \(\mathbf{K}\) 的一个纯超越扩张。*
\end{enumerate}

\BourbakiExerciseGroup

\begin{enumerate}
\def\labelenumi{\arabic{enumi})}
\tightlist
\item
  设 \(\mathbf{K}\) 是特征为 \(p > 0\) 的域，\(\mathbf{E}\) 是 \(\mathbf{K}\) 的一个扩张，\(\mathbf{B}\) 是一个 \(p\) -基（\(\mathbf{E}\) 在 \(\mathbf{K}(\mathbf{E}^p)\) 上的）（V，第 98 页）。
\end{enumerate}

\begin{enumerate}
\def\labelenumi{\alph{enumi})}
\item
  假设 \(\mathbf{E} \subset \mathbf{K}^{p^{-1}}\) ；设 \(\mathbf{K}_0\) 是 \(\mathbf{K}\) 的一个子域，使得 \(\mathbf{E}\) 在 \(\mathbf{K}_0\) 上可分。证明集合 \(\mathbf{B}^p\) 是一个 \(p\) -无关子集（在 \(\mathbf{K}_0(\mathbf{K}^p)\) 上）（注意，若 \((a_{\lambda})\) 是向量 \(K_{0}\) -空间 K 的一组基，则 \((a_{\lambda}^{p})\) 是 \(K_{0}(K^{p})\) 在 \(K_{0}\) 上的一组基。设 C 是 K 的一个子集，与 \(B^{p}\) 不相交，且使 \(B^{p} \cup C\) 是 K 在 \(K_{0}(K^{p})\) 上的一个 p-基；证明 \(B \cup C\) 是 E 在 \(K_{0}(E^{p})\) 上的一个 p-基。
\item
  假设 \(E \subset K^{p^{-n}}\) ，且 E 在 K 的一个子域 \(K_{0}\) 上可分。证明：若 K 在 \(K_{0}\) 上的不完全度有限（V，第 170 页，习题 1），则它等于 E 在 \(K_{0}\) 上的不完全度（利用 a））。
\end{enumerate}

\begin{enumerate}
\def\labelenumi{\arabic{enumi})}
\setcounter{enumi}{1}
\tightlist
\item
  设 \(\mathbf{K}\) 是特征为 \(p > 0\) 的域，\(\mathbf{E}\) 是 \(\mathbf{K}\) 的一个扩张，\(\mathbf{F}\) 是 \(\mathbf{E}\) 的一个可分扩张。证明下列三个性质中的任意两个蕴含第三个：
\end{enumerate}

\begin{enumerate}
\def\labelenumi{\alph{enumi})}
\item
  B 是一个 \(p\) -基（E 在 \(\mathbf{K}(\mathbf{E}^p)\) 上的）；
\item
  C 是一个 \(p\) -基（\(\mathbf{F}\) 在 \(\operatorname{E}(\mathbf{F}^p)\) 上的）；
\item
  \(\mathbf{B} \subset \mathbf{E}\) ，\(\mathbf{B} \cup \mathbf{C}\) 是一个 \(p\) -基（\(\mathbf{F}\) 在 \(\mathbf{K}(\mathbf{F}^p)\) 上的），且 \(\mathbf{B} \cap \mathbf{C} = \emptyset\) 。（利用如下事实：若 \((c_{\mu})\) 是向量 E-空间 \(\mathbf{F}\) 的一组基，则 \((c_{\mu}^p)\) 是 \(\mathbf{K}(\mathbf{E}^p)[\mathbf{F}^p]\) 在 \(\mathbf{K}(\mathbf{E}^p)\) 上以及 \(\mathbf{E}[\mathbf{F}^p]\) 在 \(\mathbf{E}\) 上的一组基。）
\end{enumerate}

\begin{enumerate}
\def\labelenumi{\arabic{enumi})}
\setcounter{enumi}{2}
\item
  设 K 是特征为 p \textgreater{} 0 的域，E 是 K 的一个扩张，F 是 E 的一个有限生成扩张。证明 F 在 K 上的不完全度至少等于 E 在 K 上的不完全度（归结为以下两种情形： \(1^{\circ}\) F 在 E 上可分； \(2^{\circ}\) F = E(x)，其中 \(x^{p} \in E\) （参见 V，第 170 页，习题 2））。
\item
  设 \(\mathbf{F}\) 是域 \(\mathbf{K}\) —— 特征为 \(p > 0\) —— 的一个可分扩张。证明：若 \(\mathbf{E}\) 是 \(\mathbf{F}\) 的一个在 \(\mathbf{K}\) 上相对完全的子扩张（V，第 170 页，习题 1），则 \(\mathbf{F}\) 在 \(\mathbf{E}\) 上可分。
\item
  设 \(\mathbf{K}\) 是特征为 \(p > 0\) 的域。证明：若 \(\mathbf{E}\) 是 \(\mathbf{K}\) 的代数扩张或 \(\mathbf{K}\) 的相对完全扩张，则 \(\mathrm{E}^{p^{-\infty}} = \mathrm{E}(\mathrm{K}^{p^{-\infty}})\) 。
\item
  设 \(\mathbf{K}\) 是特征为 \(p > 0\) 的域，\(\mathbf{E}\) 是 \(\mathbf{K}\) 的一个可分扩张，\(\mathbf{B}\) 是一个 \(p\) -基（\(\mathbf{E}\) 在 \(\mathbf{K}(\mathbf{E}^p)\) 上的）。
\end{enumerate}

\begin{enumerate}
\def\labelenumi{\alph{enumi})}
\item
  证明 B 在 K 上是代数自由的（考虑 B 的元素之间次数最低的一个代数关系，并将其中出现的变元的次数写成 \(kp + h\) 的形式，其中 \(0 \leqslant h \leqslant p - 1\) ）。由此推出：若 tr. \(\deg_{K}E < +\infty\) ，则 E 在 K 上的不完全度至多等于 tr. \(\deg_{K}E\) 。
\item
  证明 E 在 K(B) 上是可分的且相对完美的。
\end{enumerate}

\begin{enumerate}
\def\labelenumi{\arabic{enumi})}
\setcounter{enumi}{6}
\item
  证明特征为 p \textgreater{} 0 的域 K 的一个相对完全超越扩张 E 在 K 上不是有限生成的。（在相反的情形，证明对 E 在 K 上的每个超越基 B，E 都在 K(B) 上是代数的且可分的。）
\item
  若 E 是特征为 p \textgreater{} 0 的域 K 的一个可分代数扩张，则 E 的最大完全子扩张 \(E^{p^{\infty}}\) （诸域 \(E^{p^{n}}\) （ \(n \geqslant 0\) ）的交）在 K 的最大完全子域 \(K^{p^{\infty}}\) 上是代数的。（若 \(x \in E^{p^{\infty}}\) ，证明对每个整数 n \textgreater{} 0，我们有 \(x^{p^{-n}} \in \mathrm{K}(x)\) ，并由此推出 x 在 \(K^{p^{n}}\) 上的极小多项式与它在 K 上的极小多项式相同。）
\item
  设 \(\mathbf{E}\) 是一个域，\(\mathbf{G}\) 是 \(\mathbf{E}\) 的一个自同构群，而 \(\mathbf{K}\) 是 \(\mathbf{E}\) 的由在 \(\mathbf{G}\) 下不变的元素组成的子域。
\end{enumerate}

\begin{enumerate}
\def\labelenumi{\alph{enumi})}
\item
  证明：对 E 的每个在 G 的元素下稳定的子域 L，L 与 K 在 L ∩ K 上线性不相交。（考虑一个线性关系 \(\sum_{j}\lambda_{j}x_{j}=0\) ，其中的元素 \(x_{j} \in \mathbf{K}\) 在 \(\mathbf{L} \cap \mathbf{K}\) 上线性无关，而 \(\lambda_{j} \in \mathbf{L}\) 不全为零，且 \(\lambda_{j} \neq 0\) 的个数尽可能地少。）
\item
  若 \(L \cap K\) 是 K 的某个自同构群的不变域，证明 L 是 \(\mathrm{L}(\mathbf{K})\) 的某个自同构群的不变域。
\end{enumerate}

\begin{enumerate}
\def\labelenumi{\arabic{enumi})}
\setcounter{enumi}{9}
\item
  设 E 是域 K 的一个扩张，G 是拓扑群 \(\operatorname{Aut}_{\mathbf{K}}(\mathbf{E})\) 的一个紧子群，F 是 G 的不变域。证明 E 是 F 的一个伽罗瓦扩张，且 G 典范地同构于拓扑群 \(\operatorname{Aut}_{\mathbf{F}}(\mathbf{E})\) 。（注意 G 以连续方式作用于离散空间 E，以确立元素 \(x \in E\) 在 G 的作用下具有有限轨道。）
\item
  设 \(\mathbf{K}\) 是特征为 \(p > 0\) 的域。对每个交换 K-代数 \(\mathbf{A}\) ，取 \(p\) 次幂使环 \(\mathbf{A}\) 成为一个 A-代数，从而也是一个 K-代数，记作 \(\mathbf{A}^{p^{-1}}\) 。设 \(\mathbf{A}\) 是一个交换 K-代数。证明下列条件是等价的：
\end{enumerate}

\begin{enumerate}
\def\labelenumi{(\roman{enumi})}
\item
  \(\mathbf{A}\) 是一个可分的 K-代数。
\item
  唯一的 K-同态 \(\Phi: \mathbf{A} \otimes_{\mathbf{K}} \mathbf{K}^{p^{-1}} \to \mathbf{A}^{p^{-1}}\) 使得 \(\Phi(a \otimes \lambda) = a^p \cdot \lambda (a \in \mathbf{A}, \lambda \in \mathbf{K})\) 是单射的。
\item
  对每个族 \((b_{i})_{i\in I}\) （在 \(\mathbf{K}\) 中、于 \(\mathbf{K}^p\) 上线性无关）以及每个族 \((a_{i})_{i\in I}\) （在 \(\mathbf{A}\) 中），等式 \(\sum_{i}a_i^p b_i = 0\) 蕴含 \(a_{i} = 0\) 对所有 \(i\) 成立。
\end{enumerate}

（可以验证（ii） \(\Leftrightarrow\) （iii），并且（ii）是V，第123页，定理2，e）的推论。）

\begin{enumerate}
\def\labelenumi{\arabic{enumi})}
\setcounter{enumi}{11}
\tightlist
\item
  设 K 是一个域，\((\mathbf{X}_{i})_{i\in I}\) 是一族未定元。证明形式幂级数代数 \(\mathbf{K}[[(\mathbf{X}_{i})_{i\in I}]]\) 是一个可分 K-代数。（应用习题 11；也可以验证，对 K 的每个有限扩张 \(K'\) ，\(\mathbf{K}[[(\mathbf{X}_{i})_{i\in I}]]\otimes_{\mathbf{K}}K'\) 同构于 \(\mathbf{K}'[[(\mathbf{X}_{i})_{i\in I}]]\) ，然后应用 V，第 123 页，定理 2，d)。）由此推出 \(\mathbf{K}[[(\mathbf{X}_{i})_{i\in I}]]\) 的分式域是 K 的一个可分扩张（V，第 120 页，命题 4）。
\end{enumerate}

\BourbakiExerciseGroup

\begin{enumerate}
\def\labelenumi{\arabic{enumi})}
\tightlist
\item
  设 \(\mathbf{K}\) 是特征为 \(p > 0\) 的域，E 是 \(\mathbf{K}\) 的一个可分扩张，且在 \(\mathbf{K}\) 上具有有限超越次数。
\end{enumerate}

\begin{enumerate}
\def\labelenumi{\alph{enumi})}
\item
  若存在 E 在 K 上的一个超越基 \(B_{0}\) 以及一个整数 \(m \geqslant 0\) ，使得 \(\mathrm{K}(\mathrm{E}^{p^{m}})\) 在 \(\mathrm{K}(\mathrm{B}_{0})\) 上可分，证明：对 E 在 K 上的每个超越基 B，存在整数 \(n \geqslant 0\) 使得 \(\mathrm{K}(\mathrm{E}^{p^{n}})\) 在 \(\mathrm{K}(\mathrm{B})\) 上可分。
\item
  由此推出：若 a) 中的条件成立，则 E 在 K 上具有一个可分超越基（若 S 是 E 在 \(\mathrm{K}(\mathrm{E}^{p})\) 上的一个 p-基，而 B 是 E 在 K 上的包含 S 的一个超越基（V，第 176 页，习题 6），则利用 V，第 176 页，习题 6，证明 B 是 E 在 K 上的一个可分超越基）。
\item
  设 K 是特征为 p \textgreater{} 0 的域，x 是 K 上的一个超越元（在 K 的一个代数闭扩张 \(\Omega\) 中）。证明诸扩张 \(\mathrm{K}(x^{p^{-n}})\) 的并 E 是 K 的一个可分扩张，满足 tr. \(\deg_{K}E = 1\) ，但不具有可分超越基。
\item
  设 K 是一个完全域。证明有理分式域 \(\mathrm{K}(\mathrm{T})\) 的完全闭包是 K 的一个超越次数为 1 的扩张，它不具有可分超越基。
\end{enumerate}

\begin{enumerate}
\def\labelenumi{\arabic{enumi})}
\setcounter{enumi}{1}
\item
  设 K 是特征为 p \textgreater{} 0 的域，E 是 K 的一个有限超越次数的扩张。要使 E 在 K 上具有一个可分超越基，必须而且只需 E 在 K 上可分，且 E 在 K 上的不完全度等于它在 K 上的超越次数。
\item
  设 \(\mathbf{E}\) 是域 \(\mathbf{K}\) 的一个扩张，\(\mathbf{F}\) 是 \(\mathbf{E}\) 的一个扩张。
\end{enumerate}

\begin{enumerate}
\def\labelenumi{\alph{enumi})}
\item
  若 E 在 K 上具有一个可分超越基，且 F 在 E 上具有一个可分超越基，则 F 在 K 上具有一个可分超越基。
\item
  若 F 在 K 上具有一个可分超越基，且 tr. \(\deg_{K}E<+\infty\) ，则 E 在 K 上具有一个可分超越基（归结为 tr. \(\deg_{K}F<+\infty\) 的情形，并应用习题 1）。
\item
  若 F 在 K 上具有一个有限可分超越基，且 F 在 E 上可分，则 F 在 E 上具有一个可分超越基（利用习题 1）。
\end{enumerate}

\begin{enumerate}
\def\labelenumi{\arabic{enumi})}
\setcounter{enumi}{3}
\item
  设 K 是特征为 p \textgreater{} 0 的域，其绝对不完全度（V，第 170 页，习题 1）等于 1。证明：要使 K 的一个扩张 E 在 K 上可分，必须而且只需 E 不包含任何在 K 上 p-根式且不属于 K 的元素（若 \(a \in K\) 构成 K 在 \(K^{p}\) 上的一个 p-基，证明 \(\mathbf{K}^{p^{-1}} = \mathbf{K}(a^{p^{-1}})\) 与 E 在 K 上线性不相交）。
\item
  \begin{enumerate}
  \def\labelenumii{\alph{enumii})}
  \tightlist
  \item
    设 K 是特征为 p \textgreater{} 0 的域，E 是 K 的一个具有可分超越基的扩张。证明诸域 \(\mathrm{K}(E^{p^{n}})\) （n 遍历 N）的交 L 是 K 的一个代数扩张。（首先考虑 \(\operatorname{tr} \cdot \deg_{K} E < +\infty\) 的情形；证明 L 在 K 上相对完全，从而 E 在 L 上可分。然后利用 V，第 176 页，习题 2 和习题 6。在一般情形，若 B 是 E 在 K 上的一个可分超越基，则对每个 \(x \in L\) ，存在 B 的一个有限子集 \(B_{0}\) ，使得 x 在每个 \(\mathrm{K}(B_{0}^{p^{n}})\) 上的极小多项式都与它在 \(\mathrm{K(B)}\) 上的极小多项式相同；由此推出，若 \(\mathrm{F} = \mathrm{K}(B_{0})(x)\) ，则对所有 n 有 \(x \in \mathrm{K}(F^{p^{n}})\) 。）
  \end{enumerate}
\end{enumerate}

\begin{enumerate}
\def\labelenumi{\alph{enumi})}
\setcounter{enumi}{1}
\tightlist
\item
  证明包含在 E 中的最大完全域 \(E^{p^{\infty}}\) （诸 \(E^{p^{n}}\) 的交）是 \(K^{p^{\infty}}\) 的一个代数扩张（利用 a) 以及 V，第 176 页，习题 8）。给出一个例子：E 是完全域 K 的有限生成扩张，但 \(E^{p^{\infty}}\) 不包含在 K 中。
\end{enumerate}

\begin{enumerate}
\def\labelenumi{\arabic{enumi})}
\setcounter{enumi}{5}
\tightlist
\item
  设 \(\mathbf{K}\) 是特征为 \(p > 0\) 的域，\(\mathbf{E}\) 是 \(\mathbf{K}\) 的一个扩张，使得存在整数 \(h \geqslant 0\) ，\(\mathbf{K}(\mathrm{E}^{p^h})\) 在 \(\mathbf{K}\) 上可分（当 \(\mathbf{E}\) 是 \(\mathbf{K}\) 的有限生成扩张时，情形总是如此）。
\end{enumerate}

\begin{enumerate}
\def\labelenumi{\alph{enumi})}
\item
  若 C 是一个 \(p\) -基（E 在 \(\mathbf{K}(\mathbf{E}^p)\) 上的），则存在 C 的一个子集 B，使得 \(\mathbf{B}^{p^h}\) 是一个 \(p\) -基（\(\mathbf{K}(\mathbf{E}^{p^h})\) 在 \(\mathbf{K}(\mathbf{E}^{p^{h+1}})\) 上的）；B 是 K 上的一个代数自由集。证明扩张 \(\mathbf{F} = \mathbf{K}(\mathbf{E}^{p^h})(\mathbf{B})\) 在 K 上可分，且 E 是 F 的代数扩张。
\item
  证明 B 是一个 \(p\) -基（F 在 \(\mathbf{K}(\mathbf{F}^p)\) 上的）。证明：若 \(x \in \mathbf{E}\) ，且 \(m\) 是使 \(x^{p^m} \in \mathbf{F}\) 的最小整数，则 \(x^{p^m} \in \mathbf{K}(\mathbf{F}^{p^m})\) ；因此 E 是域 \(\mathbf{K}^{p^{-\infty}}(\mathbf{F})\) （同构于 \(\mathbf{K}^{p^{-\infty}} \otimes_{\mathbf{K}} \mathbf{K}\) ）的一个子域。称 E 的一个子扩张 F 为可区分的，如果它在 K 上可分，且 E 是 \(\mathbf{K}^{p^{-\infty}}(\mathbf{F})\) 的一个子域；此时 F 是 K 在 E 中的一个极大可分扩张。
\item
  设 P 是特征为 p \textgreater{} 0 的完全域，\(K = P(X, Y)\) 是两个未定元的有理分式域，而 \(\Omega\) 是有理分式域 \(\mathbf{K}(\mathbf{Z})\) 在 K 上的一个代数闭扩张。若 \(u \in \Omega\) 满足 \(u^{p} = X + YZ^{p}\) ，证明 \(\mathrm{E} = \mathrm{K}(\mathrm{Z}, u)\) 是 K 的一个扩张，使得 \(\mathrm{K}(\mathrm{Z})\) 与 \(\mathrm{K}(u)\) 是 E 的两个可区分子扩张，从而 E 不存在更大的可分子扩张。若 \(v \in \Omega\) 满足 \(v^{p^{2}} = X + YZ^{p}\) ，证明在扩张 \(\mathrm{E}' = \mathrm{K}(\mathrm{Z}, v)\) 中，\(\mathrm{K}(\mathrm{Z})\) 是 E 的一个极大可分子扩张，但不是可区分子扩张。
\end{enumerate}

\begin{enumerate}
\def\labelenumi{\arabic{enumi})}
\setcounter{enumi}{6}
\tightlist
\item
  假设习题 6 的假设成立，并且 E 在 K 上的不完全度（V，第 170 页，习题 1）有限（若 E 是 K 的有限生成扩张，情形总是如此）。
\end{enumerate}

\begin{enumerate}
\def\labelenumi{\alph{enumi})}
\tightlist
\item
  设 B 是一个 \(p\) -基（F 在 \(\mathbf{K}(\mathbf{F}^p)\) 上的），它同时是 F 在 K 上的一个可分超越基（并且是 E 在 K 上的一个超越基）。证明：对每个整数 \(k > 0\) ，\(\mathbf{B}^{p^k}\) 是一个 \(p\) -无关子集（在 \(\mathbf{K}(\mathbf{E}^{p^k})\) 中、于 \(\mathbf{K}(\mathbf{E}^{p^{k+1}})\) 上）。设 \(q = \operatorname{tr} \cdot \deg_{\mathbf{K}} \mathbf{E}\) ；若 \(m_k + q\) 是 \(\mathbf{K}(\mathbf{E}^{p^k})\) 在 \(\mathbf{K}(\mathbf{E}^{p^{k+1}})\) 上的不完全度，证明 \(p^f\) —— 其中 \(f = \sum_{k=0}^{h-1} m_k\) —— 等于次数 {[}E : F{]} ，因此与
\end{enumerate}

E 的可区分子扩张 F 的选择无关。整数 \(p^{f}\) 称为 E 在 K 上的不可分度，记作 \([E:K]_{i}\) ；当 E 为有限次时，这个数与 V，第 46 页中记作 \([E:K]_{i}\) 的不可分次数一致。

\begin{enumerate}
\def\labelenumi{\alph{enumi})}
\setcounter{enumi}{1}
\item
  对 \(0 \leqslant k \leqslant h - 1\) ，\([\mathbf{K}^{p^{-k}}(\mathbf{E}) : \mathbf{K}^{p^{-k}}]_i\) 等于 \(p^{f_k}\) ，其中 \(f_k = \sum_{j=k}^{h-1} m_j\) 。
\item
  设 \(L\) 是 \(E\) 的一个子扩张，在 \(K\) 上可分，且使得 \(E\) 是 \(p\) -根式的且在 \(L\) 上有限次；证明 \([E:L] \geqslant [E:K]_i\) （若 \([L(E^{p^k}):K(E^{p^k})] = p^{r_k}\) ，证明 \(r_{k+1} \leqslant r_k + q\) ，这通过考虑 \(L(E^{p^k})\) 在 \(K(E^{p^k})\) 上的一组基以及 \(L\) 在 \(K(L^p)\) 上的一组基来完成）。
\item
  设 \(K'\) 是 K 的一个扩张，使得 \(K'\) 与 E 包含在 K 的同一个代数闭扩张 \(\Omega\) 中。证明：若 \(\mathrm{K}'(\mathrm{F})\) 在 \(K'\) 上可分，则它是 \(\mathrm{K}'(\mathrm{E})\) （ \(K'\) 的扩张）的一个可区分子扩张。当 \(K'\) 与 E 在 K 上代数不相交时即是如此，且此时我们有 \([\mathrm{K}'(\mathrm{E}):\mathrm{K}']_{i}\leqslant[\mathrm{E}:\mathrm{K}]_{i}\) ；当 \(K'\) 与 E 在 K 上线性不相交，或当 \(K'\) 是 K 的可分代数扩张时，等号成立。e) 证明 \([\mathrm{E}:\mathrm{K}]_{i}\) 是对 E 在 K 上的所有超越基 B 取的 \([\mathrm{E}:\mathrm{K}(\mathrm{B})]_{i}\) 的最小值；若 \(\bar{K}\) 是 K 在 \(\Omega\) 中的相对代数闭包（它是 K 的一个代数闭包），则 \([\mathrm{E}:\mathrm{K}(\mathrm{B})]_{i}=[\mathrm{E}:\mathrm{K}]_{i}[\bar{\mathrm{K}}(\mathrm{E}):\bar{\mathrm{K}}(\mathrm{B})]_{i}\) 。若 L 是 E 的一个子扩张，使得 E 是 L 的代数扩张，且 \(\bar{\mathrm{K}}(\mathrm{E})\) 是 \(\bar{\mathrm{K}}(\mathrm{L})\) 的（代数）可分扩张，则 L 在 E 中包含的最大可分扩张 F 是可区分的。
\item
  假设 E 是 K 的一个有限生成扩张；证明：对 E 的每个子扩张 L，我们有 \([L:K]_{i} \leqslant [E:K]_{i} \leqslant [L:K]_{i}[E:L]_{i}\) 。给出一个例子，使得
\end{enumerate}

\[
[ \mathrm{L}: \mathrm{K} ] _ {i} = [ \mathrm{E}: \mathrm{K} ] _ {i} \quad \text { and } \quad [ \mathrm{E}: \mathrm{L} ] _ {i} > 1.
\]

\begin{enumerate}
\def\labelenumi{\arabic{enumi})}
\setcounter{enumi}{7}
\tightlist
\item
  设 \(k\) 是一个域，\(L\) 是 \(k\) 的一个扩张。证明下列条件是等价的：
\end{enumerate}

\begin{enumerate}
\def\labelenumi{(\roman{enumi})}
\item
  L 是一个有限生成的纯超越扩张的可分代数扩张。
\item
  \(L\) 是 \(k\) 的一个可分扩张，且
\end{enumerate}

\[
[ \Omega_ {L / k}: L ] = \operatorname{tr} \cdot \deg_ {k} L <   + \infty .
\]

（为证明 (ii) \(\Rightarrow\) (i)，我们可以如下进行：设 L 中的 \(x_{1},\ldots ,x_{n}\) 使得 \(dx_{1},\ldots ,dx_{n}\) 构成 \(\Omega_k(\mathbf{L})\) 的一组基，并令 \(\mathbf{K} = k(x_1,\dots,x_n)\) 。首先证明 \(\operatorname {tr.deg}_k(\mathbf{K}) = n\) ，从而 K 是 \(k\) 的一个纯超越扩张。设 \(\mathbf{L}'\) 是包含在 L 中的 K 的一个有限代数扩张。通过将典范映射 \(u:\Omega_k(\mathbf{K})\otimes_\mathbf{K}\mathbf{L}'\to \Omega_k(\mathbf{L}')\) 与典范映射 \(\Omega_k(\mathbf{K})\otimes_\mathbf{K}\mathbf{L}\to \Omega_k(\mathbf{L})\) 相比较，证明 \(u\) 是单射。计算 \([\Omega_k(\mathbf{L}'):\mathbf{L}']\) （V，第 134 页，推论 1）。由此推出 \(u\) 是双射，从而 \(\Omega_{\mathbf{K}}(\mathbf{L}') = 0\) 。证明 \(\mathbf{L}'\) 在 K 上可分；由此推出 L 在 K 上可分。）

\begin{enumerate}
\def\labelenumi{\arabic{enumi})}
\setcounter{enumi}{8}
\tightlist
\item
  沿用习题 8 的记号，设 \(\mathbf{M} \subset \mathbf{L}\) 是一个子扩张。令 \(m = \operatorname{tr} \cdot \deg_k \mathbf{M}\) ，\(n = \operatorname{tr} \cdot \deg_k \mathbf{L}\) ，\(r = n - m\) 。设 \(x_1, \ldots, x_n\) （在 \(\mathbf{L}\) 中）使得 \(\mathbf{L}\) 是 \(k(x_1, \ldots, x_n) = \mathbf{K}\) 的一个可分代数扩张。
\end{enumerate}

\begin{enumerate}
\def\labelenumi{\alph{enumi})}
\item
  证明：若有必要，在将 \(x_{i}\) 重新编号后，我们可以假设 \(L\) 在 \(M(x_{1},\ldots ,x_{r})\) 上是代数的。
\item
  证明存在一个有限生成的子扩张 \(\mathbf{N} \subset \mathbf{M}\) ，使得：
\end{enumerate}

\(\alpha)\) L 在 \(\mathbf{N}(x_1,\dots,x_s)\) 上是代数的；

\(\beta)\) \(\mathbf{K}' = \mathbf{K}(\mathbf{N}(x_1,\dots,x_r))\) 是 \(\mathbf{N}(x_1,\dots,x_r)\) 的一个有限扩张；

\(\gamma)\) 典范同态 \(\mathbf{K}^{\prime}\otimes_{\mathbf{N}(x_1,\ldots ,x_r)}\mathbf{M}(x_1,\dots,x_r)\to \mathbf{K}(\mathbf{M}(x_1,\dots,x_r)) = \mathbf{K}''\) 是双射的。

\begin{enumerate}
\def\labelenumi{\alph{enumi})}
\setcounter{enumi}{2}
\tightlist
\item
  设 N 为 b) 中那样的一个子扩张。依次证明：\(K''\) 是 \(K'\) 的可分代数扩张（因为 L 在 K 上可分）；\(M(x_{1}, \ldots, x_{r})\) 是 \(N(x_{1}, \ldots, x_{r})\) 的可分代数扩张（利用 b)，\(\gamma)\)）；以及 M 是 N 的可分代数扩张。d) 证明 M 是一个纯超越扩张的可分代数扩张（利用 c) 以及 N 是 k 的有限生成可分扩张这一事实）。
\end{enumerate}

\BourbakiExerciseGroup

\begin{enumerate}
\def\labelenumi{\arabic{enumi})}
\tightlist
\item
  设 \(K_{0}\) 是特征为 p \textgreater{} 0 的域，\(\mathbf{K} = \mathbf{K}_{0}(\mathbf{X}, \mathbf{Y})\) 是 \(K_{0}\) 上两个未定元的有理分式域，U 与 V 是两个未定元，\(\Omega\) 是 \(\mathbf{K}(\mathbf{U}, \mathbf{V})\) 的一个代数闭扩张，\(\mathbf{E} = \mathbf{K}(\mathbf{U}, u)\) ，其中 \(u^{p} = X + YU^{p}\) ，而 \(F = K(V, v)\) ，其中 \(v^{p} = X + YV^{p}\) ；E 与 F 在 K 上不是可分的（V，第 178 页，习题 6，c））。
\end{enumerate}

\begin{enumerate}
\def\labelenumi{\alph{enumi})}
\item
  证明 \(\mathbf{K}\) 在 \(\mathbf{E}\) 中与在 \(\mathbf{F}\) 中都是相对代数封闭的（若 \(x \in \mathbf{E}\) 在 \(\mathbf{K}\) 上是代数的，证明必有 \(x^p \in \mathbf{K}\) ；若 \(x \notin \mathbf{K}\) ，则 \(\mathbf{E} = \mathbf{K}(\mathbf{U}, x)\) ，由此推出 \(\mathbf{X}^{1/p}\) 与 \(\mathbf{Y}^{1/p}\) 将属于 \(\mathbf{K}(x)\) ）。
\item
  证明 E 与 F 在 K 上线性不相交，但 K 在 \(\mathbf{K}(\mathbf{E} \cup \mathbf{F})\) 中不是相对代数封闭的。（证明 \(X^{1/p} \in \mathbf{K}(\mathbf{E} \cup \mathbf{F})\) ；由此推出 v 不能属于 \(\mathbf{E}(\mathbf{V})\) ，并由此得出 E 与 F 在 K 上线性不相交。）
\end{enumerate}

\begin{enumerate}
\def\labelenumi{\arabic{enumi})}
\setcounter{enumi}{1}
\tightlist
\item
  设 \(\mathbf{K}\) 是特征为 \(p > 0\) 的域，\(\Omega\) 是 \(\mathbf{K}\) 的一个扩张，而 \(\mathbf{E}\) 、\(\mathbf{F}\) 是 \(\Omega\) 的两个在 \(\mathbf{K}\) 上代数不相交的子扩张。
\end{enumerate}

\begin{enumerate}
\def\labelenumi{\alph{enumi})}
\tightlist
\item
  设 \(F_{s}\) 是 K 在 F 中的相对可分闭包。证明 \(E(F_{s})\) 是 E 在 \(E(F)\) 中的相对可分闭包。（归结为 \(F_{s}=K\) 的情形；设 B 是 E 在 K 上的一个超越基，而 \(K'\) 是 K 在 F 中的相对代数闭包，它在 F 上是根式的；则（V，第 137 页）\(K'(B)\) 是 \(K(B)\) 在 \(F(B)\) 中的相对代数闭包，且在 \(K(B)\) 上是 p-根式的。若 \(x\in E(F)\) 在 E 上是代数的，则存在整数 \(r\geqslant0\) 使得 \(x^{p'}\in M\) ，其中 M 是 \(F(B)\) 的一个有限次可分扩张，具有一组基 \((u_{i})_{1\leqslant i\leqslant n}\) （在 \(\mathrm{F(B)}\) 上，由 E 的元素组成）。若 \(y = \sum_{i}b_{i}u_{i}\) ，其中 \(b_{i}\in \mathrm{F(B)}\) ，证明
\end{enumerate}

通过考虑 M 到 E(F) 的一个代数闭包中的 F(B)-同构，证明 \(b_{i}\) 属于 K'(B)，并由此推出 y 在 E 上是 p-根式的。）

\begin{enumerate}
\def\labelenumi{\alph{enumi})}
\setcounter{enumi}{1}
\tightlist
\item
  假设 E 与 F 在 K 上线性不相交，K 在 F 中相对代数封闭，且 E 是 K 的可分扩张。证明 E 在 E(F) 中是相对代数封闭的。（借助 V，第 138 页，归结为 E 是 K 的可分代数扩张的情形。由 a) 推出 E 在 E(F) 中的相对代数闭包在 E 上是 p-根式的，然后利用如下事实：若 \((a_{\lambda})\) 是 E 在 K 上的一组基，则诸 \(a_{\lambda}^{p'}\) 也构成 E 在 K 上的一组基。）
\end{enumerate}

\begin{enumerate}
\def\labelenumi{\arabic{enumi})}
\setcounter{enumi}{2}
\tightlist
\item
  设 E 是一个域，G 是 E 的一个自同构群，而 K 是 G 的不变域。a) 证明：若 G 中不存在异于 G 本身的有限指数正规子群，则 E 是 K 的正则扩张。
\end{enumerate}

\begin{enumerate}
\def\labelenumi{\alph{enumi})}
\setcounter{enumi}{1}
\tightlist
\item
  由 a) 推出：若 G 同构于无限单群的一个有限直积，则 E 是 K 的正则扩张。
\end{enumerate}

\begin{enumerate}
\def\labelenumi{\arabic{enumi})}
\setcounter{enumi}{3}
\tightlist
\item
  设 K 为一个域，E、F 为 K 的两个扩张。
\end{enumerate}

\begin{enumerate}
\def\labelenumi{\alph{enumi})}
\item
  假设 K 在 F 中相对代数封闭。证明：若 \((L, u, v)\) 与 \((L', u', v')\) 是 E 与 F 的两个复合扩张，使得 \(u(E)\) 与 \(v(F)\) （相应地 \(u'(E)\) 与 \(v'(F)\) ）在 K 上代数不相交，则这两个复合扩张是同构的。（首先考虑 E 是 K 的纯超越扩张的情形；然后利用习题 2，a)，以及 V，第 154 页，习题 6，b)）
\item
  设 \(E_{s}\) 、\(F_{s}\) 分别是 K 在 E 与 F 中的相对可分闭包。假设 \(E_{s}\) 与 \(F_{s}\) 的所有复合扩张都是同构的。证明：若 \((L, u, v)\) 与 \((L', u', v')\) 是 E 与 F 的两个复合扩张，使得 \(u(E)\) 与 \(v(F)\) （相应地 \(u'(E)\) 与 \(v'(F)\) ）在 K 上代数不相交，则这两个复合扩张是同构的（利用 a) 以及 V，第 173 页，习题 11）。
\end{enumerate}

\begin{enumerate}
\def\labelenumi{\arabic{enumi})}
\setcounter{enumi}{4}
\tightlist
\item
  设 K 是一个域，p 是它的特征指数。对每个与 p 互素的整数 n \textgreater{} 0 ，用 \(\varphi_{\mathrm{K}}(n)\) 表示 \(\mathrm{R}_{n}(\mathrm{K}) = \mathrm{K}(\mu_{n})\) 在 K 上的次数，其中 \(\mu_{n}\) 表示 K 的一个代数闭包中 n 次单位根所成的群。例如 \(\varphi_{\mathrm{Q}}(n) = \varphi(n)\) （V，第 84 页，定理 2）；若 K 是含 q 个元素的有限域，则 \(\varphi_{\mathrm{K}}(n)\) 是 q 的剩余类在 \((\mathbf{Z}/n\mathbf{Z})^{*}\) 中的阶（V，第 97 页，推论）；若 K 是代数闭的，则对所有 n 有 \(\varphi_{\mathrm{K}}(n) = 1\) 。
\end{enumerate}

\begin{enumerate}
\def\labelenumi{\alph{enumi})}
\item
  设 E 是素域 \(K_{0}\) 在 K 中的相对代数闭包。证明 \(\varphi_{\mathrm{K}}(n) = \varphi_{\mathrm{E}}(n)\) 。
\item
  假设 K 是 \(K_{0}\) 的有限生成扩张。设 m 是整数 \textgreater0 ；则存在整数 N \textgreater0 ，它被每个满足 \(\varphi_{\mathrm{K}}(n) \leqslant m\) 的 n 整除。
\item
  设 \(s\) 是 \(\mathrm{GL}_m(\mathbf{K})\) 的一个有限阶元素。证明 \(s^N\) 的阶是 \(p\) 的一个幂。（我们有 \(s = tu = ut\) ，其中 \(u\) 的阶是 \(p\) 的幂，而 \(t\) 的阶 \(n\) 与 \(p\) 互素。元素 \(t\) 是可分多项式 \(\mathbf{X}^n - 1\) 的根，同时也是其特征多项式的根；由此推出子代数 \(\mathbf{A}\) —— \(\mathbf{M}_m(\mathbf{K})\) 的、由 \(t\) 生成的 —— 是平展的，次数 \(\leqslant m\) ，同构于一些扩张 \(\mathbf{R}_d(\mathbf{K})\) 的积，其中 \(d\) 遍历 \(n\) 的某些因子，包括 \(n\) 本身。因此 \(\varphi_{\mathbf{K}}(n) \leqslant m\) ；利用 \(b\) ）。）
\end{enumerate}
